\section{Weighted distribution modules and exact rational-grid ranks}
\label{sec:distribution}

The distribution relations in the Stieltjes drafts have a simple
representation-theoretic structure. The relevant dimension is the
dimension of a formal quotient by specified relations. The classical
theory of universal distributions provides the broader setting
\cite{Kubert,Ouyang}. Here we give an explicit character normal form,
including its behavior at exceptional weights and in finite Taylor
jets. This proves the two all-denominator rank statements suggested by
the source computations, without making an independence assertion about
evaluated special values.

\subsection{The relation module and its complete character decomposition}

Fix \(q\geq2\), and put
\[
 G_q=\frac1q\Z/\Z,\qquad U(q)=(\Z/q\Z)^\times.
\]
Let \(\mathcal A\) be a commutative unital \(\C\)-algebra and choose
arbitrary \(t_p\in\mathcal A\), one for each prime \(p\mid q\).
The choices may vanish or be nilpotent. Let
\[
 \mathcal V_q=\bigoplus_{x\in G_q}\mathcal A[x].
\]
For \(p\mid q\) and \(x\in pG_q\), define the prime distribution row
\begin{equation}
 D_p(x)=\sum_{\substack{y\in G_q\\py=x}}[y]-t_p[x].
 \label{eq:dist-prime-row}
\end{equation}
Write \(\mathcal R_q\) for the \(\mathcal A\)-span of these rows and
\(\mathcal U_q=\mathcal V_q/\mathcal R_q\).
The unit group acts by \(u[x]=[ux]\), and each relation map is
equivariant.

For \(d\mid q\), set \(t_d=\prod_{p\mid d}t_p^{v_p(d)}\).
Rows for composite divisors introduce no additional relations. Indeed,
if \(ab\mid q\) and \(x\in abG_q\), then
\begin{equation}
 D_{ab}(x)=
 \sum_{\substack{z\in G_q\\az=x}}D_b(z)+t_bD_a(x).
 \label{eq:dist-composite-row}
\end{equation}
Every \(z\) in this sum belongs to \(bG_q\), so the rows on the right
are legitimate grid rows. Expanding the two sums proves the identity.
Induction on the number of prime factors of \(d\), counted with
multiplicity, proves the assertion for every divisor.

Let \(\chi\) run through all characters of \(U(q)\), and let \(f=f_\chi\)
be the conductor of its primitive inducing character \(\chi^*\).
For \(f\mid d\mid q\), let \(\chi_d\) denote the induced character
modulo \(d\), and define
\begin{equation}
 Y_{d,\chi}=\sum_{a\in U(d)}\chi_d(a)[a/d].
 \label{eq:dist-character-vector}
\end{equation}
For the principal character at \(d=1\), this means
\(Y_{1,\mathbf1}=[0]\). These are formal vectors; no numerical
\(L\)-value is being declared independent.

\begin{lemma}[All character components of the rows]
\label{lem:dist-all-components}
The \(\chi^{-1}\)-eigenspace of \(\mathcal V_q\) is free over
\(\mathcal A\), with basis
\[
 \{Y_{d,\chi}: f_\chi\mid d\mid q\}.
\]
Its relation submodule is generated exactly by the edges
\begin{equation}
 \begin{aligned}
 Y_{pd,\chi}-t_pY_{d,\chi},
     &\qquad p\mid d,\\
 Y_{pd,\chi}-(t_p-\chi^*(p))Y_{d,\chi},
     &\qquad p\nmid d,
 \end{aligned}
 \qquad f_\chi\mid d,\quad pd\mid q.
 \label{eq:dist-character-edges}
\end{equation}
In particular, there are no additional character relations arising
from denominators smaller than, or not divisible by, \(f_\chi\).
\end{lemma}

\begin{proof}
The elements of \(G_q\) of exact denominator \(d\) form a single
\(U(q)\)-orbit, naturally identified with \(U(d)\). Reduction
\(U(q)\to U(d)\) is surjective. The character \(\chi\) factors through
this map exactly when \(f_\chi\mid d\); otherwise its projection on
that orbit is zero. More explicitly, the idempotent
\[
 E_\chi=\frac1{\varphi(q)}\sum_{u\in U(q)}\chi(u)u
\]
satisfies, for \(a\in U(d)\),
\begin{equation}
 E_\chi[a/d]=
 \begin{cases}
 \displaystyle\frac{\chi_d(a)^{-1}}{\varphi(d)}Y_{d,\chi},
       &f_\chi\mid d,\\[6pt]
 0,&f_\chi\nmid d.
 \end{cases}
 \label{eq:dist-projector}
\end{equation}
This follows by grouping \(u\) according to its residue modulo \(d\)
and using character orthogonality on the kernel of reduction.
The idempotents are mutually orthogonal and sum to the identity.
All their scalar denominators are units in \(\mathcal A\), so the
argument applies equally to a nonreduced algebra.

For a fixed \(p\mid q\), consider the equivariant map
\[
 D_p:\mathcal A[pG_q]\longrightarrow\mathcal V_q,\qquad
 [x]\longmapsto D_p(x).
\]
The exact-denominator orbits in its domain have \(pd\mid q\).
By \eqref{eq:dist-projector}, their only nonzero \(\chi\)-components
have \(f_\chi\mid d\), and each such component has the single
generator \(Y_{d,\chi}\). Consequently, computing
\(D_p(Y_{d,\chi})\) computes \emph{every} projected relation for
that \(p\); an orbit on which the projector vanishes cannot create
a hidden relation in the target.

If \(p\mid d\), all \(p\) preimages of each \(a/d\), \(a\in U(d)\),
have exact denominator \(pd\). Summing with weights \(\chi_d(a)\)
therefore gives \(Y_{pd,\chi}-t_pY_{d,\chi}\).
If \(p\nmid d\), one preimage has denominator \(d\), while the other
\(p-1\) have denominator \(pd\). In the former preimage,
\(pc\equiv a\pmod d\), so its character coefficient is
\(\chi_d(a)=\chi^*(p)\chi_d(c)\). Thus the sum is
\[
 Y_{pd,\chi}+\chi^*(p)Y_{d,\chi}-t_pY_{d,\chi}.
 \]
Here \(p\nmid f_\chi\), as \(f_\chi\mid d\), so \(\chi^*(p)\)
is a nonzero root of unity. This proves both the edge formulas
and their completeness.
\end{proof}

\begin{theorem}[Polynomial normal form]
\label{thm:dist-normal-form}
For \(f=f_\chi\mid d\mid q\), define
\begin{equation}
 A_{d,\chi}(\mathbf t)=
 \prod_{p\mid f}t_p^{\,v_p(d)-v_p(f)}
 \prod_{\substack{p\mid d\\p\nmid f}}
       t_p^{\,v_p(d)-1}(t_p-\chi^*(p)).
 \label{eq:dist-normal-coefficient}
\end{equation}
Empty products are \(1\), so \(A_{f,\chi}=1\).
Then the full \(\chi\)-relation submodule is
\begin{equation}
 \mathcal R_{q,\chi}
 =\bigoplus_{\substack{f\mid d\mid q\\d\ne f}}
   \mathcal A\bigl(Y_{d,\chi}-A_{d,\chi}Y_{f,\chi}\bigr).
 \label{eq:dist-triangular-relations}
\end{equation}
In particular,
\begin{equation}
 \mathcal U_q\cong
 \bigoplus_{\chi\in\widehat{U(q)}}\mathcal A\,Y_{f_\chi,\chi}
 \cong\mathcal A^{\varphi(q)}.
 \label{eq:dist-free}
\end{equation}
The character normal form of every vector is obtained by replacing
\(Y_{d,\chi}\) with \(A_{d,\chi}Y_{f_\chi,\chi}\).
\end{theorem}

\begin{proof}
The polynomials \(A_{d,\chi}\) satisfy precisely the edge recursions
\eqref{eq:dist-character-edges}: a prime already dividing \(d\)
contributes \(t_p\), and a prime introduced for the first time
outside \(f\) contributes \(t_p-\chi^*(p)\).
The factors for distinct primes commute.

Follow any chain of prime extensions from \(f\) to \(d\).
Successively eliminating the intermediate \(Y\)'s expresses
\(Y_{d,\chi}-A_{d,\chi}Y_{f,\chi}\) as a combination of the
edge relations. Conversely, write an edge as
\(Y_{pd,\chi}-cY_{d,\chi}\), where
\(A_{pd,\chi}=cA_{d,\chi}\). It equals
\[
 \bigl(Y_{pd,\chi}-A_{pd,\chi}Y_{f,\chi}\bigr)
 -c\bigl(Y_{d,\chi}-A_{d,\chi}Y_{f,\chi}\bigr).
 \]
Lemma~\ref{lem:dist-all-components} therefore gives equality of the
two relation spans.

The displayed generators in \eqref{eq:dist-triangular-relations}
are independent: the coefficient of \(Y_{d,\chi}\), for \(d\ne f\),
in a linear combination is exactly the coefficient of its own
generator. Together with \(Y_{f,\chi}\), they form a triangular
basis of the whole character block. Its quotient is consequently
free of rank one. Summing over all \(\varphi(q)\) characters proves
\eqref{eq:dist-free}.
\end{proof}

\begin{remark}
The proof divides by neither \(t_p\) nor \(t_p-\chi^*(p)\).
It proves freeness first over \(\C[t_p:p\mid q]\) and then over
every \(\C\)-algebra obtained by specialization. Exceptional weights
can therefore make a preferred coordinate vanish, but cannot change
the rank of the full quotient. This distinction is essential when
Euler factors are used to solve a relation system.
\end{remark}

\subsection{Anchors, reflection, and the ranks in the drafts}

In this subsection take \(\mathcal A=\C\).
In a Hurwitz-zeta grid, use the representatives \(0<a\leq1\):
the group element \(0\) represents \(a=1\), not the singular
Hurwitz parameter \(a=0\). Treating the endpoint as known means
imposing the homogeneous anchor \([0]=0\). It removes exactly the
principal generator \(Y_{1,\mathbf1}\). Hence
\begin{equation}
 \dim_\C\mathcal V_q/(\mathcal R_q+\C[0])=\varphi(q)-1.
 \label{eq:dist-anchor-dimension}
\end{equation}
There are \(q-1\) remaining grid variables, so their distribution
coefficient matrix has rank \(q-\varphi(q)\). At rational weights
\(t_p\), the matrix has rational entries and this is also its rank
over \(\Q\).

Let \(J[x]=[-x]\). On the \(\chi\)-block, \(J\) acts as
\(\chi(-1)\). Thus the additional homogeneous reflection rows
\begin{equation}
 [x]+(-1)^k[-x]=0,\qquad x\ne0,
 \label{eq:dist-reflection-rows}
\end{equation}
retain exactly the characters with
\(\chi(-1)=(-1)^{k+1}\). After anchoring, the omission of \(x=0\)
in \eqref{eq:dist-reflection-rows} has no effect.

\begin{corollary}[All-denominator formal ranks]
\label{cor:dist-all-ranks}
For \(q\geq3\) and every integer \(k\geq0\), the formal quotient
by distribution, endpoint, and the reflection rows
\eqref{eq:dist-reflection-rows} has dimension
\begin{equation}
 r_{q,k}=
 \begin{cases}
 \varphi(q)/2-1,&k\ \text{odd},\\
 \varphi(q)/2,&k\ \text{even}.
 \end{cases}
 \label{eq:dist-parity-dimension}
\end{equation}
This conclusion holds for arbitrary prime weights.
In particular it proves the all-\(q\), all-\(k\) formal rank law
for the negative Hurwitz-zeta derivative ladder.
The formal distribution-only rank in the Stieltjes parameter
derivative tower is \(q-\varphi(q)\) for every \(q\geq2\),
every parameter derivative order \(k\geq1\), and every Stieltjes
index \(n\geq1\), relative to the endpoint and lower-index data.
\end{corollary}

\begin{proof}
For \(q\geq3\), the element \(-1\) is a nonidentity involution in
\(U(q)\). Character orthogonality gives
\(\sum_\chi\chi(-1)=0\), and each summand is \(1\) or \(-1\).
There are consequently \(\varphi(q)/2\) characters of each parity.
Reflection retains the parity stated above. The principal
character is even, so the endpoint removes one additional
dimension exactly when \(k\) is odd. This proves
\eqref{eq:dist-parity-dimension}.

For the negative ladder, differentiation of the Hurwitz
multiplication formula at \(s=-k\) gives
\begin{equation}
 \sum_{j=0}^{d-1}
 \zeta'\!\left(-k,\frac{a+j}{d}\right)
 =d^{-k}\bigl(\zeta'(-k,a)+\log d\,\zeta(-k,a)\bigr).
 \label{eq:dist-negative-jet}
\end{equation}
The Bernoulli-polynomial values \(\zeta(-k,a)\) are lower data.
Thus the homogeneous distribution weight is \(t_p=p^{-k}\).
The homogeneous reflection is
\eqref{eq:dist-reflection-rows}; an explicit right-hand side is
given below. This applies the first part to the ladder in the
source article.

For the parameter derivative tower, the Laurent expansion is
\begin{equation}
 \zeta(s,a)=\frac1{s-1}
   +\sum_{n\geq0}\frac{(-1)^n}{n!}\gamma_n(a)(s-1)^n.
 \label{eq:dist-stieltjes-laurent}
\end{equation}
For \(k\geq1\), its pole disappears after differentiation in \(a\):
\begin{equation}
 F_k(\epsilon,a):=
 \partial_a^k\zeta(1+\epsilon,a)
 =(-1)^k(1+\epsilon)_k\zeta(k+1+\epsilon,a)
 =\sum_{n\geq0}\frac{(-1)^n}{n!}
          \gamma_n^{(k)}(a)\epsilon^n.
 \label{eq:dist-stieltjes-generating}
\end{equation}
This standard identity is the starting point of the all-order
parameter derivative formula of Coffey \cite{Coffey}.
Differentiating the Hurwitz multiplication formula \(k\) times
in \(a\), including its chain-rule factors, gives
\[
 \sum_{j=0}^{d-1}
 F_k\!\left(\epsilon,\frac{a+j}{d}\right)
 =d^{k+1+\epsilon}F_k(\epsilon,a).
\]
Comparing coefficients yields the exact triangular identity
\begin{equation}
 \sum_{j=0}^{d-1}
 \gamma_n^{(k)}\!\left(\frac{a+j}{d}\right)
 =d^{k+1}\sum_{r=0}^{n}
   \binom nr(-\log d)^{n-r}\gamma_r^{(k)}(a).
 \label{eq:dist-stieltjes-multiplication}
\end{equation}
With lower indices \(r<n\) treated as specified data, the
homogeneous weight is \(t_p=p^{k+1}\), independently of \(n\).
Equation~\eqref{eq:dist-anchor-dimension} now gives rank
\(q-\varphi(q)\).
\end{proof}

For completeness, the right-hand side of the negative reflection
can be written without introducing another Hurwitz derivative.
For \(0<a<1\), set
\[
 C_j(a)=\sum_{m\geq1}\frac{\cos(2\pi ma)}{m^j},
 \qquad
 S_j(a)=\sum_{m\geq1}\frac{\sin(2\pi ma)}{m^j}.
\]
For \(k\geq1\), these series at \(j=k+1\) converge absolutely, and
the Fourier form of the Hurwitz functional equation \cite{DLMF}
gives
\begin{equation}
 \zeta'(-k,a)+(-1)^k\zeta'(-k,1-a)
 =\frac{(-1)^{\lfloor k/2\rfloor}k!}{(2\pi)^k}
 \begin{cases}
 C_{k+1}(a),&k\ \text{even},\\
 S_{k+1}(a),&k\ \text{odd}.
 \end{cases}
 \label{eq:dist-negative-reflection}
\end{equation}
Indeed, the Fourier expansion of \(\zeta(s,a)\) has the factor
\[
 \frac{2\Gamma(1-s)}{(2\pi)^{1-s}}
 \left(\sin\frac{\pi s}{2}\,C_{1-s}(a)
       +\cos\frac{\pi s}{2}\,S_{1-s}(a)\right).
\]
The indicated reflection selects a sine or cosine prefactor
that vanishes at \(s=-k\). Only the derivative of that prefactor
survives, giving \eqref{eq:dist-negative-reflection}.
This also explains why reflection is an inhomogeneous relation
with depth-one data in the inventory.

\begin{corollary}[The formal gamma grid]
\label{cor:dist-gamma-rank}
For \(q\geq3\), the \(q-1\) formal coordinates
\(\log\Gamma(a/q)\), \(1\leq a<q\), modulo multiplication and
reflection and relative to their stated elementary right-hand
sides, have quotient dimension \(\varphi(q)/2\).
\end{corollary}

\begin{proof}
Lerch's identity
\(\zeta'(0,a)=\log\Gamma(a)-\tfrac12\log(2\pi)\)
translates gamma multiplication into
\eqref{eq:dist-negative-jet} with \(k=0\), hence \(t_p=1\).
The reflection formula
\[
 \zeta'(0,a)+\zeta'(0,1-a)=-\log(2\sin\pi a)
\]
has homogeneous part \([x]+[-x]=0\). Apply
\eqref{eq:dist-parity-dimension} with \(k=0\).
\end{proof}

\begin{remark}[What these rank theorems settle]
\label{rem:dist-formal-only}
The first Stieltjes source states a finite verification followed
by an all-denominator conjecture; the parameter-derivative source
states a general rank theorem while reporting only a finite
calculation. Corollary~\ref{cor:dist-all-ranks} supplies a proof
for the specified relation modules in both cases. The README
had already identified the second proof gap.
Corollary~\ref{cor:dist-gamma-rank} likewise determines the rank
of the stated gamma relation matrix. None of these results says
that the evaluated numbers admit no further relations.
In particular, the claim that multiplication and reflection
generate every algebraic multiplicative relation among rational
gamma values remains part of the conjectural completeness
problem, not a consequence of the matrix rank.
\end{remark}

\subsection{Analytic weights, exceptional coordinates, and finite jets}

Specialize \(t_p=p^s=\exp(s\log p)\), using the real logarithm
of the prime. Formula~\eqref{eq:dist-normal-coefficient} becomes
\begin{equation}
 A_{d,\chi}(s)=
 \left(\frac d{f_\chi}\right)^s
 \prod_{\substack{p\mid d\\p\nmid f_\chi}}
       \left(1-\chi^*(p)p^{-s}\right).
 \label{eq:dist-euler-coefficient}
\end{equation}
This is exactly the familiar imprimitive Euler-factor pattern
for Dirichlet \(L\)-functions \cite{DLMF}. The normal-form proof
shows why it appears in the formal relations before any
evaluation is made.

The symbols of exact denominator \(q\) have character coordinates
\(Y_{q,\chi}\). In the quotient, their \(\chi\)-coordinate is
\begin{equation}
 Y_{q,\chi}=A_{q,\chi}(s)Y_{f_\chi,\chi}.
 \label{eq:dist-top-coordinate}
\end{equation}
They form a basis of the full quotient precisely when every
\(A_{q,\chi}(s)\) is nonzero. In particular this holds when
\(\Repart s\ne0\), since \(\chi^*(p)\) has modulus one.
All nonzero integer weights in the two Stieltjes articles are
therefore free of this coordinate obstruction. On the imaginary
axis the top-denominator coordinates can fail, while the
conductor coordinates of Theorem~\ref{thm:dist-normal-form}
remain a basis.

\begin{theorem}[Exact defect of top-denominator jets]
\label{thm:dist-jet-defect}
Fix \(s_0\in\C\) and \(N\geq0\), and work over
\(\mathcal A_N=\C[\epsilon]/(\epsilon^{N+1})\) with
\(t_p=p^{s_0+\epsilon}\), truncated in this ring.
The full distribution module is free of rank \(\varphi(q)\)
over \(\mathcal A_N\), hence has complex dimension
\((N+1)\varphi(q)\).
Set
\begin{equation}
 r_\chi(s_0)=
 \#\{p\mid q:p\nmid f_\chi,\ p^{s_0}=\chi^*(p)\}.
 \label{eq:dist-resonance-count}
\end{equation}
The image of the top-denominator submodule in its \(\chi\)-block
has complex dimension
\begin{equation}
 \max\{N+1-r_\chi(s_0),0\},
 \label{eq:dist-jet-image}
\end{equation}
and its cokernel has dimension
\(\min\{r_\chi(s_0),N+1\}\).
An endpoint anchor or a fixed-parity reflection restricts these
formulas to the retained characters. The resulting full quotient
remains free over \(\mathcal A_N\) with the ranks already computed.
\end{theorem}

\begin{proof}
Freeness is Theorem~\ref{thm:dist-normal-form}, applied to
\(\mathcal A_N\). A resonant factor in
\eqref{eq:dist-euler-coefficient} has a simple zero at \(s_0\):
\[
 \left.\frac{d}{ds}
 (1-\chi^*(p)p^{-s})\right|_{s=s_0}=\log p\ne0.
\]
All other factors, including \((q/f_\chi)^s\), are nonzero
there. In the ring of analytic germs,
\[
 A_{q,\chi}(s_0+\epsilon)
   =\epsilon^{r_\chi(s_0)}u_\chi(\epsilon),
 \qquad u_\chi(0)\ne0.
\]
After truncation, \(u_\chi\) is a unit. The
top-denominator image in the free rank-one block is therefore
the ideal \(\epsilon^{r_\chi(s_0)}\mathcal A_N\).
Its basis consists of the powers
\(\epsilon^{r_\chi},\ldots,\epsilon^N\) when \(r_\chi\leq N\),
and it is zero otherwise. This proves both dimensions.
Anchoring and reflection remove entire character blocks,
so the same calculation applies to each retained block.
\end{proof}

Comparing Taylor coefficients of a weighted distribution identity
gives the corresponding stacked finite-jet system. If one identifies
the coefficient at order \(j\) with
\(\epsilon^{N-j}[x]\), its relation space is the
\(\mathcal A_N\)-span above; factorial normalizations merely
rescale coordinates. Thus freeness also gives the rank of the
complete simultaneous jet system, rather than just its successive
highest layers.

\begin{proposition}[Multiplicity of nonzero resonances]
\label{prop:dist-resonance-multiplicity}
Resonances occur only on \(\Repart s_0=0\).
For \(s_0\ne0\), each \(r_\chi(s_0)\) is at most one.
Multiple resonant primes in a single character block can occur
only at \(s_0=0\).
\end{proposition}

\begin{proof}
Taking absolute values in \(p^{s_0}=\chi^*(p)\) gives
\(\Repart s_0=0\). Write \(s_0=\ii t\), with \(t\ne0\).
Because every character value is a root of unity, resonance
at \(p\) implies
\(t\log p/(2\pi)\in\Q\). This rational number is nonzero.
If distinct primes \(p\) and \(\ell\) both resonate, then
\(\log p/\log\ell\in\Q\). Writing this positive rational as
\(a/b\) would give \(p^b=\ell^a\) for positive integers \(a,b\),
contradicting unique factorization.
\end{proof}

For a concrete origin resonance, take \(q=14\) and let \(\chi^*\)
be the quadratic character modulo \(7\). Its conductor is \(7\),
its parity is odd, and \(\chi^*(2)=1\). Then
\begin{equation}
 Y_{14,\chi}(s)=(2^s-1)Y_{7,\chi}(s).
 \label{eq:dist-resonance-example}
\end{equation}
At \(s=0\), the top-denominator coordinate vanishes although the
conductor direction persists. For any analytic realization of
this relation,
\begin{align}
 Y_{14,\chi}'(0)&=\log2\,Y_{7,\chi}(0),\notag\\
 Y_{14,\chi}''(0)&=
 2\log2\,Y_{7,\chi}'(0)+(\log2)^2Y_{7,\chi}(0).
 \label{eq:dist-resonance-example-jets}
\end{align}
Thus recovering the first conductor derivative from these
top-denominator data requires the second top-denominator jet.
The actual realization
\[
 Y_{d,\chi}(s)
 =\sum_{a\in U(d)}\chi_d(a)\zeta(s,a/d)
 =d^sL(s,\chi_d)
\]
obeys the same identity. This example concerns an imprimitive
Euler factor; it is distinct from the trivial-zero regularization
of the primitive \(L\)-function in the next section.
